\chapter{Introduction}

\section{Background}

When doing image analysis in the spatial domain, operations are performed directly
on pixels. Blurring an image, involves sliding an averaging kernel over
the pixel grid; edge detection applies a Sobel operator to the same grid. But,
spatial filtering is inherently limited. Applying a blur to suppress noise in a complex scene
will degrade the structures that should be preserved: the edges, textures, and
boundaries of the objects of interest.

\vspace{0.15in}
\noindent
The frequency domain offers a fundamentally different perspective \cite{oppenheim1999}.
A large, uniformly coloured region such as a clear sky corresponds to low-frequency
content. A sharp edge such as the silhouette of a mountain against that sky corresponds
to high-frequency content. Random grain or sensor noise manifests as very
high-frequency energy \cite{gonzalez_woods}. By operating in this domain, one gains the
ability to precisely target specific types of content for removal or enhancement without
disturbing the rest of the image \cite{pratt2007}.

\vspace{0.15in}
\noindent
Our project, which we are calling as \textbf{di3p}, is built to actually implement this idea in
C++ without leaning on OpenCV or any other library that would just hide what was
happening underneath. We wanted to write the transforms ourselves and see exactly what was going on at every step.

\section{Motivation}

Existing frequency-domain filtering techniques like Butterworth and Gaussian 
filters are applied globally across an entire image \cite{gonzalez_woods}. 
This uniform treatment ignores the fact that distinct regions of a natural image 
exhibit different spectral characteristics \cite{pratt2007}. A flat sky 
and a textured flower petal will require different denoising strategies. 
While image segmentation methods such as K-Means \cite{jain_kmeans_survey} and 
Mean Shift \cite{comaniciu2002} are well-studied, they are rarely composed with 
per-region frequency-domain processing into a unified pipeline. So now, our project solves 
that gap by combining K-Means spatial segmentation with region-specific DFT and DCT 
filtering.

\vspace{0.15in}
\noindent
Consider photo of a flower with a distracting repeating-pattern background. A spatial blur will smooth both the noise and the flower simultaneously, destroying
structural detail in the attempt to remove unwanted content. A frequency-domain approach
allows us to identify the spectral signature of the noise and suppress it selectively.

\vspace{0.15in}
\noindent
First, we performed performed segmentation-first processing. If the flower and
the background are separated into distinct clusters before any frequency-domain work is
performed, filters can be tailored per-region. The flower's cluster might receive a
sharpening high-pass filter to emphasise petal edges, and the background cluster
receives a low-pass filter to soften the noise pattern. This chaining of K-Means spatial
grouping with DFT/DCT frequency analysis is the main idea behind out project.
